% =============================================================================
% CAPÍTULO 8 — VERIFICACIÓN, VALIDACIÓN Y ENTORNOS DE TRABAJO
% Archivo maestro del capítulo. Importa las secciones desde sections/cap08/.
% =============================================================================
\chapter{Verificación, Validación y Entornos de Trabajo}
\label{ch:vyv}
\resumenCapitulo{Este capítulo cubre cómo se comprueba y se ejecuta el software: la
estrategia de pruebas (Testcontainers en el backend, Vitest en el frontend), la
configuración de los entornos de desarrollo y producción, y la gestión de configuraciones
y secretos conforme a la norma ISO/IEC/IEEE~15289:2024.}

La norma ISO/IEC/IEEE~15289:2024 exige documentar no solo \textit{qué} hace el software,
sino \textbf{cómo se comprueba que lo hace correctamente} y \textbf{en qué entornos se
ejecuta}. Este capítulo cubre los tres ítems de información correspondientes: la estrategia
de pruebas (verificación y validación), la configuración de los entornos de ejecución
(desarrollo y producción) y la gestión de configuraciones y secretos.

Estos aspectos se apoyan en el diseño descrito en capítulos anteriores: las pruebas
verifican el comportamiento del backend (capítulo~\ref{ch:backend},
pág.~\pageref{ch:backend}) y del frontend (capítulo~\ref{ch:frontend},
pág.~\pageref{ch:frontend}); los entornos materializan el empaquetado monolítico
(sección~\ref{subsec:empaquetado}, pág.~\pageref{subsec:empaquetado}); y la gestión de
secretos extiende el modelo de seguridad de las integraciones externas
(capítulo~\ref{ch:integraciones}, pág.~\pageref{ch:integraciones}).

\begin{NotaTecnica}
La distinción entre \textbf{verificación} (``¿construimos el producto correctamente?'',
es decir, conforme a su especificación) y \textbf{validación} (``¿construimos el producto
correcto?'', es decir, el que satisface la necesidad) guía la organización de este
capítulo: las pruebas automatizadas verifican; los entornos y la gestión de configuración
soportan la validación en condiciones realistas.
\end{NotaTecnica}

% --- Secciones del capítulo (orden del índice instrucciones.md §8) ---
% =============================================================================
% 8.1 ESTRATEGIA DE PRUEBAS DE SOFTWARE
% =============================================================================
\section{Estrategia de Pruebas de Software}
\label{sec:pruebas}

La estrategia de pruebas de EthosHub se reparte entre los dos componentes del sistema con
herramientas idiomáticas a cada \textit{stack}: \textbf{Testcontainers} y \textbf{JUnit}
en el backend, \textbf{Vitest} y \textbf{Testing Library} en el frontend. La
figura~\ref{fig:piramide} sitúa cada nivel de prueba en la clásica \textbf{pirámide de
pruebas}.

\vspace{0.3cm}
\begin{figure}[h!]
\centering
\begin{tikzpicture}[font=\sffamily\scriptsize]
    \fill[colorAcento!30, draw=colorAcento!80!black] (0,3) -- (-1.1,1.8) -- (1.1,1.8) -- cycle;
    \fill[c4Contenedor!40, draw=c4Sistema!70!black] (-1.1,1.8) -- (-2.3,0.6) -- (2.3,0.6) -- (1.1,1.8) -- cycle;
    \fill[c4Componente!45, draw=c4Sistema!70!black] (-2.3,0.6) -- (-3.5,-0.6) -- (3.5,-0.6) -- (2.3,0.6) -- cycle;
    \node at (0,2.35) {E2E};
    \node[align=center] at (0,1.15) {Integración\\\scriptsize Testcontainers + JUnit};
    \node[align=center] at (0,-0.05) {Unitarias / Renderizado\\\scriptsize Vitest · Testing Library · JUnit};
    \draw[-{Stealth[length=2mm]}, grisISO] (4.4,-0.6) -- node[right, text=grisISO, font=\sffamily\tiny]{realismo} (4.4,3.0);
    \draw[-{Stealth[length=2mm]}, grisISO] (-4.4,3.0) -- node[left, text=grisISO, font=\sffamily\tiny]{volumen} (-4.4,-0.6);
\end{tikzpicture}
\caption{Pirámide de pruebas de EthosHub por nivel y herramienta.}
\label{fig:piramide}
\end{figure}

\subsection{Pruebas de integración del backend con Testcontainers}
\label{subsec:testcontainers}

Las pruebas de integración del backend ejecutan contra un \textbf{PostgreSQL real},
levantado en un contenedor efímero mediante \textbf{Testcontainers}. Así se elimina la
brecha entre el entorno de prueba y el de producción: el SQL tipado de jOOQ
(capítulo~\ref{ch:backend}, sección~\ref{subsec:jooq}, pág.~\pageref{subsec:jooq}) se
verifica contra el mismo motor que se usará en producción, no contra una base de datos
en memoria con dialecto distinto.

\clearpage
\begin{lstlisting}[language=Java, caption={Configuración de Testcontainers con \texttt{@ServiceConnection} (\texttt{TestcontainersConfiguration.java}).}, label={lst:testcontainers}]
@TestConfiguration(proxyBeanMethods = false)
class TestcontainersConfiguration {
    @Bean
    @ServiceConnection
    PostgreSQLContainer<?> postgresContainer() {
        return new PostgreSQLContainer<>(
            DockerImageName.parse("postgres:latest"));
    }
}
\end{lstlisting}

La anotación \comando{@ServiceConnection} (Spring Boot) conecta automáticamente el
contenedor con el \textit{datasource} de la aplicación, sin configuración manual de URL,
usuario o contraseña. La tabla~\ref{tab:tests-backend} enumera las clases de prueba reales
del backend y su objetivo.

\vspace{0.2cm}
\noindent
\renewcommand{\arraystretch}{1.3}
\fontsize{8.5}{10.2}\selectfont\sffamily
\begin{tabularx}{\textwidth}{|L{5.0cm}|Y|}
\hline
\rowcolor{headerTabla}
\sffamily\textbf{Clase de prueba} & \sffamily\textbf{Verifica} \\
\hline
\texttt{EthosHubApplicationTests} & Que el contexto de Spring arranca correctamente. \\ \hline
\rowcolor{grisTecnico}
\texttt{AuthControllerSecurityTest} & La protección de los \textit{endpoints} de autenticación. \\ \hline
\texttt{AuthServiceOAuthSyncTest} & La sincronización de perfil tras un \textit{login} OAuth. \\ \hline
\rowcolor{grisTecnico}
\texttt{UserRoleTest} & La derivación y validación del rol de usuario. \\ \hline
\end{tabularx}
\label{tab:tests-backend}

\begin{NotaTecnica}
La pila de pruebas del backend (declarada en \texttt{pom.xml}) combina
\texttt{spring-boot-starter-test}, \texttt{spring-boot-testcontainers},
\texttt{testcontainers-junit-jupiter}, \texttt{testcontainers-postgresql} y
\texttt{spring-security-test}. Esta última habilita la simulación de usuarios autenticados
con roles concretos, esencial para verificar la matriz de permisos del
capítulo~\ref{ch:requisitos} (pág.~\pageref{ch:requisitos}).
\end{NotaTecnica}

\subsection{Pruebas del frontend con Vitest y Testing Library}
\label{subsec:vitest_vyv}

El frontend usa \textbf{Vitest 4} sobre un entorno virtual \textbf{jsdom}, con
\textbf{Testing Library} para interactuar con los componentes como lo haría una persona
usuaria (consultas por rol accesible, eventos de usuario), no por detalles de
implementación. La configuración vive en \pathbreak{vite.config.ts} (bloque \texttt{test})
y carga \pathbreak{src/test/setup.ts}, como resume la tabla~\ref{tab:vitest}.

\vspace{0.2cm}
\noindent
\renewcommand{\arraystretch}{1.3}
\fontsize{8.5}{10.2}\selectfont\sffamily
\begin{tabularx}{\textwidth}{|L{3.4cm}|Y|}
\hline
\rowcolor{headerTabla}
\sffamily\textbf{Parámetro} & \sffamily\textbf{Valor} \\
\hline
\texttt{environment} & \texttt{jsdom} (DOM virtual en Node). \\ \hline
\rowcolor{grisTecnico}
\texttt{globals} & \texttt{true} (API de pruebas global sin importaciones). \\ \hline
\texttt{setupFiles} & \pathbreak{./src/test/setup.ts} (matchers de jest-dom). \\ \hline
\rowcolor{grisTecnico}
\texttt{alias} & \texttt{@ $\rightarrow$ ./src} (coherente con el \textit{build}). \\ \hline
\end{tabularx}
\label{tab:vitest}

El ejemplo canónico es \pathbreak{src/test/authFlow.test.tsx}, que valida el flujo de
autenticación de extremo a extremo dentro del cliente (formulario, \textit{store},
redirección por rol), apoyándose en el \textit{custom hook} \texttt{useAuthFlow}
(capítulo~\ref{ch:frontend}, sección~\ref{subsec:hooks}, pág.~\pageref{subsec:hooks}).

\begin{AvisoNota}
El \textit{linting} con \textbf{ESLint} y la verificación de tipos con \textbf{TypeScript}
estricto (capítulo~\ref{ch:frontend}, sección~\ref{sec:calidad_fe},
pág.~\pageref{sec:calidad_fe}) constituyen una primera línea de verificación
\textbf{estática} que se ejecuta antes que cualquier prueba: detectan una clase entera de
errores sin necesidad de ejecutar el código.
\end{AvisoNota}
     % 8.1 Estrategia de pruebas
% =============================================================================
% 8.2 CONFIGURACIÓN DE ENTORNOS DE EJECUCIÓN
% =============================================================================
\section{Configuración de Entornos de Ejecución}
\label{sec:entornos}

EthosHub opera en dos entornos con topologías distintas: en \textbf{desarrollo}, frontend
y backend corren como procesos separados comunicados por un \textit{proxy}; en
\textbf{producción}, ambos se fusionan en un único \textit{JAR} monolítico servido desde un
contenedor. La figura~\ref{fig:entornos} contrasta ambas topologías.

\vspace{0.3cm}
\begin{figure}[h!]
\centering
\begin{tikzpicture}[node distance=0.6cm and 1.1cm,
    box/.style={rectangle, rounded corners=3pt, draw, font=\sffamily\scriptsize, align=center,
        minimum width=2.4cm, minimum height=0.8cm, inner sep=3pt},
    grp/.style={rectangle, rounded corners=5pt, draw=grisISO!60, dashed, inner sep=8pt},
    rel/.style={-{Stealth[length=2mm]}, semithick}]
    % Desarrollo
    \node[box, fill=c4Componente!45] (vite) at (0,1.4) {Vite\\\scriptsize :3000};
    \node[box, fill=c4Datos!22, draw=c4Datos!70!black] (spring1) at (0,0) {Spring\\\scriptsize :8080};
    \draw[rel] (vite) -- node[c4etiqueta, right]{\scriptsize proxy /api} (spring1);
    \begin{scope}[on background layer]
      \node[grp, fit=(vite)(spring1), label={[font=\sffamily\scriptsize\bfseries, text=colorPrimario]above:Desarrollo}] {};
    \end{scope}
    % Produccion
    \node[box, fill=c4Sistema, text=white] (jar) at (5.4,0.7) {JAR monolítico\\\scriptsize SPA + API :8080};
    \node[box, fill=c4Externo!30, draw=c4Externo!70!black, below=of jar] (tect) at (5.4,-0.5) {+ Tectonic};
    \draw[rel] (jar) -- (tect);
    \begin{scope}[on background layer]
      \node[grp, fit=(jar)(tect), label={[font=\sffamily\scriptsize\bfseries, text=colorPrimario]above:Producción (Docker)}] {};
    \end{scope}
\end{tikzpicture}
\caption{Topologías de los entornos de desarrollo y de producción.}
\label{fig:entornos}
\end{figure}

\subsection{Entorno de desarrollo: \textit{proxy} inverso de Vite}
\label{subsec:entorno_dev}

En desarrollo, el servidor de Vite (puerto \textbf{3000}) sirve la SPA con recarga en
caliente, mientras el backend Spring corre por separado en el puerto \textbf{8080}. Para
que el cliente llame a \pathbreak{/api} sin tropezar con CORS, Vite expone un
\textbf{\textit{proxy} inverso} que reenvía las rutas \pathbreak{/api}, \pathbreak{/auth} y
\pathbreak{/v1} al backend.

\begin{lstlisting}[language=JavaScript, caption={Proxy inverso de desarrollo (\texttt{vite.config.ts}).}, label={lst:vite-proxy}]
server: {
  port: 3000,
  proxy: {
    '/api':  { target: 'http://localhost:8080', changeOrigin: true },
    '/auth': { target: 'http://localhost:8080', changeOrigin: true },
    '/v1':   { target: 'http://localhost:8080', changeOrigin: true },
  },
}
\end{lstlisting}

\begin{NotaTecnica}
El \textit{proxy} establece además \texttt{Cross-Origin-Opener-Policy:
same-origin-allow-popups} para soportar el flujo de ventana emergente de OAuth
(capítulo~\ref{ch:frontend}, sección~\ref{subsec:oauth}, pág.~\pageref{subsec:oauth}) y una
lista de \texttt{allowedHosts} para los dominios de previsualización. Este \textit{proxy}
\textbf{solo existe en desarrollo}: en producción, el mismo origen sirve la SPA y la API,
por lo que CORS deja de ser un problema.
\end{NotaTecnica}

\subsection{Entorno de producción: contenedor Docker}
\label{subsec:entorno_prod}

En producción, el sistema se empaqueta como una \textbf{imagen Docker multi-etapa}
(capítulo~\ref{ch:backend}, sección~\ref{subsec:docker}, pág.~\pageref{subsec:docker}): una
etapa de compilación con el JDK~17 y una etapa de ejecución con el JRE~17 que incorpora el
compilador \textbf{Tectonic}. El contenedor expone el puerto \textbf{8080} y arranca con
\comando{java -jar app.jar}. La tabla~\ref{tab:prod-params} resume sus parámetros.

\vspace{0.2cm}
\noindent
\renewcommand{\arraystretch}{1.3}
\fontsize{8.5}{10.2}\selectfont\sffamily
\begin{tabularx}{\textwidth}{|L{3.6cm}|Y|}
\hline
\rowcolor{headerTabla}
\sffamily\textbf{Parámetro} & \sffamily\textbf{Valor} \\
\hline
Imagen base (ejecución) & \texttt{eclipse-temurin:17-jre-jammy}. \\ \hline
\rowcolor{grisTecnico}
Puerto expuesto & \texttt{8080}. \\ \hline
Perfil Spring activo & \texttt{prod} (vía \texttt{SPRING\_PROFILES\_ACTIVE}). \\ \hline
\rowcolor{grisTecnico}
Compilador LaTeX & Tectonic 0.15.0 (caché precalentada). \\ \hline
\textit{Entrypoint} & \texttt{java -jar app.jar}. \\ \hline
\end{tabularx}
\label{tab:prod-params}

\begin{AvisoNota}
La configuración por contenedor sigue el principio de \textbf{paridad
desarrollo-producción}: el mismo artefacto que se prueba localmente (con su compilador
LaTeX y su perfil de Spring) es el que se despliega, reduciendo la clase de errores que
solo aparecen ``en producción''. Los parámetros que varían entre entornos se inyectan
exclusivamente por variables de entorno, no por reconstrucción de la imagen, como detalla
la sección~\ref{sec:secretos} (pág.~\pageref{sec:secretos}).
\end{AvisoNota}
    % 8.2 Entornos de ejecución
% =============================================================================
% 8.3 GESTIÓN DE CONFIGURACIONES Y SECRETOS
% =============================================================================
\section{Gestión de Configuraciones y Secretos}
\label{sec:secretos}

La gestión de secretos es transversal a todo el sistema y crítica para su seguridad. El
principio rector es \textbf{ningún secreto en el código fuente}: toda credencial llega por
variable de entorno y queda fuera del control de versiones. Esta sección documenta cómo se
externaliza la configuración y cómo se impide el rastreo de secretos en el repositorio.

\subsection{Propiedades de Spring y variables del entorno}
\label{subsec:propiedades}

El backend usa el mecanismo de \textbf{perfiles de Spring}
(capítulo~\ref{ch:backend}, sección~\ref{subsec:config_perfiles},
pág.~\pageref{subsec:config_perfiles}): un \pathbreak{application.yml} base con sobrecargas
\pathbreak{application-local.yml} y \pathbreak{application-prod.yml}. Los valores sensibles
se referencian con la sintaxis \comando{\$\{VARIABLE:defecto\}}, de modo que el fichero
versionado contiene el \textbf{nombre} de la variable, nunca su valor. La
tabla~\ref{tab:secretos} clasifica las credenciales por su naturaleza y ubicación.

\vspace{0.2cm}
\noindent
\renewcommand{\arraystretch}{1.3}
\fontsize{8.5}{10.2}\selectfont\sffamily
\begin{tabularx}{\textwidth}{|L{4.2cm}|C{2.0cm}|Y|}
\hline
\rowcolor{headerTabla}
\sffamily\textbf{Credencial} & \sffamily\textbf{Ámbito} & \sffamily\textbf{Ubicación} \\
\hline
\texttt{SUPABASE\_SERVICE\_ROLE\_KEY} & Privada & Solo backend (entorno). \\ \hline
\rowcolor{grisTecnico}
\texttt{GEMINI\_API\_KEY} & Privada & Solo backend (entorno). \\ \hline
\texttt{supabase.jwt-secret} & Privada & Solo backend (entorno). \\ \hline
\rowcolor{grisTecnico}
\texttt{VITE\_SUPABASE\_ANON\_KEY} & Pública (RLS) & Cliente (\textit{bundle}). \\ \hline
\texttt{VITE\_GOOGLE\_CLIENT\_ID} & Pública (origen) & Cliente (\textit{bundle}). \\ \hline
\end{tabularx}
\label{tab:secretos}

\subsection{Exclusión de secretos del control de versiones}
\label{subsec:gitignore}

El repositorio aplica una política estricta de \textbf{exclusión} mediante
\pathbreak{.gitignore}, cuyo objetivo declarado es ``no subir \textit{build}, basura, datos
de usuario, secretos ni \textit{tooling} de IA''. La figura~\ref{fig:secretos-flujo}
muestra la frontera entre lo versionado y lo excluido.

\vspace{0.3cm}
\begin{figure}[h!]
\centering
\begin{tikzpicture}[node distance=0.55cm and 1.3cm,
    ok/.style={rectangle, rounded corners=3pt, draw=c4Datos!70!black, fill=c4Datos!15,
        font=\sffamily\scriptsize, align=center, minimum width=3.0cm, minimum height=0.7cm, inner sep=3pt},
    no/.style={rectangle, rounded corners=3pt, draw=red!70!black, fill=bgAdvertencia,
        font=\sffamily\scriptsize, align=center, minimum width=3.0cm, minimum height=0.7cm, inner sep=3pt},
    rel/.style={-{Stealth[length=2mm]}, semithick}]
    \node[ok] (a) {\texttt{.env.example}\\\scriptsize (plantilla, sin valores)};
    \node[ok, below=of a] (b) {\texttt{application.yml}\\\scriptsize (solo nombres de variables)};
    \node[no, right=2.2cm of a] (c) {\texttt{.env} · \texttt{.env.*}\\\scriptsize (valores reales)};
    \node[no, below=of c] (d) {\texttt{application-prod.yml}\\\scriptsize (perfiles con secretos)};
    \node[font=\sffamily\scriptsize\bfseries, text=c4Datos!50!black] at (0,1.0) {Versionado};
    \node[font=\sffamily\scriptsize\bfseries, text=red!60!black] at (5.0,1.0) {Excluido (.gitignore)};
\end{tikzpicture}
\caption{Frontera entre artefactos versionados y secretos excluidos del repositorio.}
\label{fig:secretos-flujo}
\end{figure}

La tabla~\ref{tab:gitignore} documenta las reglas de exclusión más relevantes del
\pathbreak{.gitignore} real del backend.

\vspace{0.2cm}
\noindent
\renewcommand{\arraystretch}{1.3}
\fontsize{8.5}{10.2}\selectfont\sffamily
\begin{tabularx}{\textwidth}{|L{4.2cm}|Y|}
\hline
\rowcolor{headerTabla}
\sffamily\textbf{Patrón excluido} & \sffamily\textbf{Motivo} \\
\hline
\texttt{.env} · \texttt{.env.*} (salvo \texttt{.env.example}) & Variables con valores reales; se conserva la plantilla. \\ \hline
\rowcolor{grisTecnico}
\texttt{application-prod.yml} · \texttt{application-secret*} & Perfiles de Spring con credenciales. \\ \hline
\texttt{target/} · \texttt{*.jar} & Artefactos de \textit{build} (regenerables). \\ \hline
\rowcolor{grisTecnico}
Llaves, certificados y \textit{keystores} & Material criptográfico sensible. \\ \hline
\texttt{.claude/} · \texttt{.agents/} & \textit{Tooling} de IA, fuera del producto. \\ \hline
\end{tabularx}
\label{tab:gitignore}

\begin{AlertaSeguridad}
La regla \texttt{.env.*} con la excepción \texttt{!.env.example} es un patrón deliberado:
excluye \textbf{todos} los ficheros de entorno con valores reales pero conserva la
\textbf{plantilla} (sin secretos), de modo que un nuevo desarrollador sepa \textit{qué}
variables debe definir sin que ninguna credencial real toque el repositorio. Esta
disciplina, junto a la separación pública/privada de la tabla~\ref{tab:secretos}, cierra el
modelo de gestión de secretos de la plataforma.
\end{AlertaSeguridad}

\bigskip
Con esto se cierra la documentación de verificación, validación y entornos. El
capítulo~\ref{ch:anexos} (pág.~\pageref{ch:anexos}) recopila los anexos y el registro
documental que complementan este manual.
    % 8.3 Gestión de configuraciones y secretos
